\chapter{On Writing Well}

\section{The Book's Background}

``On Writing Well'' by William Zinsser is a classic guide to writing nonfiction that has influenced generations of writers since its first publication in 1976. Zinsser, a journalist, editor, and writing teacher, distilled decades of experience into practical principles for clear, effective writing. The book emphasizes simplicity, clarity, and the removal of clutter—principles that apply not just to writing, but to thinking itself. Zinsser argued that good writing is good thinking made visible, and that the best way to improve your writing is to cut everything that doesn't need to be there.

\nocite{zinsser:2006-on-writing-well}

\section{Typical Symptoms of Not Writing Well}

Bad non-fiction writing is easy to spot once you know what to look for. The most common symptom is clutter: unnecessary words that add nothing. Phrases like ``in order to'' instead of ``to,'' ``at this point in time'' instead of ``now,'' or ``due to the fact that'' instead of ``because'' signal that the writer hasn't cut what doesn't need to be there. Every extra word weakens the sentence.

Another symptom is passive voice. When you read ``the decision was made by management'' instead of ``management made the decision,'' you're reading weak writing. Passive voice hides who did what, makes sentences longer, and drains energy from the prose. It's a sign that the writer isn't taking responsibility for the action or doesn't want to be clear about who's acting.

Abstract language is another warning sign. Words like ``vision,'' ``strategy,'' ``transformation,'' and ``innovation'' sound impressive but say little. When a writer uses ``significant growth'' instead of ``34\% growth,'' or ``many companies'' instead of ``three companies,'' they're avoiding specificity. Abstract words create the appearance of meaning without actually providing it.

Jargon and complex sentences are also symptoms of bad writing. When a writer uses industry buzzwords or technical terms unnecessarily, they're often trying to sound smart rather than be clear. When sentences stretch across multiple lines with nested clauses and parenthetical asides, the writer has lost control of the structure. Good writing doesn't need complexity to sound authoritative.

Self-congratulation is another symptom. When writing is filled with words like ``reliable,'' ``global leading,'' and ``extraordinary,'' the writer is telling you how great they are rather than showing you. 

Overstressing history is another symptom. Consider this excerpt:

\begin{quote}
Our security businesses have long and distinguished operating histories. Several of our brands were established more than 75 years ago, and many of our brands originally created their categories:
\begin{itemize}
    \item Von Duprin, established in 1908, was awarded the first exit device patent;
    \item Schlage, established in 1920, was awarded the first patents granted for the cylindrical lock and the push button lock;
    \item LCN, established in 1926, created the first door closure.
\end{itemize}
\end{quote}

History is good, but overstressing it again and again often means the company has lost focus on what is happening now. Past milestones do not prove current advantage. Great companies acknowledge history and then show present results and present execution.

Repetition without content is another sign of bad writing. When company names, division names, or product names appear again and again without new information, the writer is filling space instead of adding meaning. Great companies use names as labels for concrete facts and actions. Weak companies repeat names as if saying them often enough will make them matter.

Finally, the absence of the writer's voice is a symptom. When writing sounds like it came from a template or a committee, when every sentence could have been written by anyone, the writer has stripped away what makes writing interesting. Writing that sounds like a manual or a legal document has lost the human element that makes non-fiction engaging.

\section{The History of the Chinese Eight-Legged Essay}

The eight-legged essay was a rigid writing format used in the Chinese imperial examination system from the Ming dynasty onward. It was designed to test candidates on their knowledge of the Confucian classics and their ability to follow a prescribed structure. The name comes from the essay's eight sections, each with a specific function and style, which candidates had to follow precisely. Mastery of this form became essential for anyone seeking a government position.

Over time, the eight-legged essay evolved from a tool for testing knowledge into an art of formula. Students memorized patterns, stock phrases, and model essays. Success depended less on original thought and more on reproducing the correct structure and approved interpretations of the classics. A talented candidate was not necessarily the one who understood the world best, but the one who could fit words into the required mold most smoothly.

By the late imperial period, criticism of the eight-legged essay was widespread. Reformers argued that it stifled independent thinking and trapped the educated elite in a closed world of textual games. Yet the system persisted because it was predictable, controllable, and familiar. Generations of scholars devoted their lives to perfecting a writing style that rewarded conformity over creativity.

When the Opium Wars and subsequent foreign invasions shook China in the nineteenth century, the limitations of this system became painfully clear. The country faced new technologies, new weapons, and new political ideas, but its educated class had been trained to excel at writing essays within a fixed pattern, not at thinking for themselves about unfamiliar problems. When real change arrived, almost no one in the official system knew how to respond. The eight-legged essay had produced experts in form, not in judgement.

\section{Eight-Legged Essays in Annual Reports}

There are actually quite a lot of eight-legged essays in annual reports. Companies fill the contents like the ancient Chinese students did: they follow prescribed formats, use stock phrases, and reproduce approved structures. Annual reports have become formulaic documents where companies slot in the required sections—executive summary, financial highlights, risk factors, corporate governance—without necessarily engaging in independent thinking about their actual business situation.

However, the business world is brutal and requires independent thinking. Markets change, technologies disrupt, competitors emerge, and customer preferences shift. Companies that simply fill in the template, that reproduce the same language year after year, that follow the formula without understanding what it means, are unlikely to survive. Like the Qing dynasty, which collapsed when it could no longer adapt to a changing world, companies that write eight-legged essays instead of thinking independently about their challenges will eventually fail. The form may look correct, but the substance is missing, and in business, substance is what matters.

\section{Business World Eight Leggy Essays}

Here are some excerpts of these eight-legged essays in annual reports.

\begin{quote}
    Together with our customers, we not only rise to meet today's challenges, but we are also shaping the future by empowering each other to continuously improve and deliver better outcomes. Technology is essential in addressing the most profound economic, environmental, and social challenges. Software and analytics help predict and prevent disasters, make cities smarter, and foster equity in workplaces.
\end{quote}
    
It talks about everything, yet says nothing. What value does the software bring to customers?

\begin{quote}
    The Group's aspirations are subject to various external and internal factors, some of which it cannot influence. Successful
achievement of the bank's 2025 strategic targets may be adversely impacted by reduced revenue generating capacities
of some of the bank's core businesses should downside risks crystallize. These risks include but are not limited to the
challenging macroeconomic environment in Europe, in particular Germany, and the potential for a re-emergence of
inflationary pressures impacting the interest rate outlook. A number of geopolitical trends may exacerbate these risks,
including potential for widespread imposition of trade tariffs by the new U.S. administration.
\end{quote}

The challenge applies to all international companies, and you forgot to mention climate change and the Ukraine war.
Come on, you don't need so many words to say it. What's the unique risk for the company?

\begin{quote}
    What a remarkable year it has been! Put simply, 2024 was outstanding for Deutsche Telekom: all of our key financial performance indicators clearly exceeded market expectations. In November, for the first time in over two decades, our share surpassed the 30 euro mark.
\end{quote}

Focusing on the scoreboard is bad; focusing on share price is worse. I understand your bonus hangs on the share price and you don't need to tell me. By the way, around 2000 Telekom's share price surpassed 100 euro, and afterwards there was a crash. So what's there to be proud of?

\begin{quote}
    We have a strong foundation for future success. Building on our position of 
strength as a leading technology company, Siemens has launched the 
ONE Tech Company program to achieve the next level of performance and value 
creation. The program aims to ensure that the company leverages the 
opportunities arising from the historic market shifts and from technological 
disruptions. 
\end{quote}

Lots of self-congratulations: strong, strength, next level. Remove these terms, and nothing is left.
